\section{Background}
\label{sec:background}

\subsection{Monolithic, modular-monolith and microservice deployments}
\label{sec:bg-styles}

A monolith packages every module of an application into one deployable unit. Modules call each
other in-process, and they usually share one database. Lewis and Fowler described the
microservice style as a suite of small services, each running in its own process, communicating
through lightweight mechanisms such as HTTP resource APIs and deployed
independently~\cite{Fowler2014}. Newman~\cite{Newman2015} and Richardson~\cite{Richardson2018}
turned this description into design guidance, including decomposition along business boundaries and
the database-per-service pattern, in which each service owns its data and other services reach it
only through the service's interface.

Between the two sits the modular monolith: a single deployment unit whose modules are separated by
enforced boundaries and explicit interfaces~\cite{SuLi2024}. A recent systematic review of modular
monoliths in cloud environments names the lack of empirical benchmarks as one of the main open
gaps~\cite{AlQoran2025}. Faustino et al.\ migrated one system step by step from a monolith to a
modular monolith and then to microservices, measuring performance and migration effort at each
step, and observed performance costs already at the modular-monolith step~\cite{Faustino2024}.



In this paper the three styles differ in one respect only: the mapping of modules to processes.
The business code is the same in every configuration (Section~\ref{sec:testbed}).

\subsection{Where distribution costs time}
\label{sec:bg-costs}

When a module call leaves its process, several costs appear that an in-process call does not pay.
Arguments and results are serialised and parsed, here as JSON. The bytes travel over TCP, and
through TLS if the connection is encrypted. Both ends run an HTTP stack. In Kubernetes, the
destination is a service virtual IP that kube-proxy rules translate to a pod address, over a pod network whose
overheads depend on the CNI plugin~\cite{Qi2021}.

A sidecar service mesh adds two proxies to every hop: the caller's outbound sidecar and the
callee's inbound sidecar. Istio's published benchmarks, measured with HTTP/1.1, 1\,kB payloads,
1{,}000 requests per second and mutual TLS, report that the client- and server-side sidecars
together add 2.65\,ms at P90 and 2.91\,ms at P99 in Istio 1.10.3~\cite{Istio110Perf}, and 0.18\,ms
and 0.25\,ms in Istio 1.22.3~\cite{Istio122Perf}. The overhead is therefore strongly
version-dependent. Zhu et al.\ found that sidecars raise end-to-end latency by 27\% to 269\% and
CPU usage by 42\% to 163\%, depending on protocol mode and enabled filters; protocol parsing
dominates in HTTP and gRPC mode, and inter-process communication and socket writes dominate in TCP
mode~\cite{Zhu2023}. Bremler-Barr et al.\ compared the P99 cost of mutual TLS across four meshes under heavy load and found
sidecar-based Istio far more expensive than Istio Ambient, Linkerd or Cilium~\cite{BremlerBarr2025}.

Distribution also changes how latency is distributed. A request that depends on several services
inherits the tails of all of them, so the tail of each component weighs on the whole~\cite{Dean2013}, and
production traces show heavy-tailed call graphs with hot-spot services~\cite{Luo2021}.

\subsection{CPU limits and CFS bandwidth control}
\label{sec:bg-cfs}

A Kubernetes CPU limit is enforced by the Linux Completely Fair Scheduler (CFS) through bandwidth
control. The container's control group receives a quota of CPU time per period, 100\,ms by
default; a limit of 500 millicores grants 50\,ms of CPU time per period, summed over all threads of
the group. When the threads have used the quota, the whole group is throttled, that is, not
scheduled, until the next period begins~\cite{KernelBWC}. The kernel exposes per-group counters of
elapsed periods, throttled periods and total throttled time~\cite{KernelBWC}, and the testbed
records them for every container (Section~\ref{sec:method}).

Throttling matters for latency more than average utilisation suggests. A multi-threaded runtime
such as the JVM, with request, garbage-collection and compiler threads, can use a small quota in a
burst early in a period. A request that arrives afterwards waits for the next period, so a
fractional CPU limit can add up to most of a 100\,ms period to that request's latency while the
container's utilisation averaged over a second stays well below its limit. Ugedal and Kumar studied
throttling in CFS bandwidth control and proposed kernel changes that reduce throttling the quota
does not require~\cite{Ugedal2022}. Splitting a fixed CPU budget across more services produces
smaller per-service limits, so decomposition interacts with this mechanism directly.

\subsection{Measuring latency under load}
\label{sec:bg-load}

Load generators follow one of two system models. In a closed system a fixed number $N$ of clients
each wait for a response, and possibly a think time, before sending the next request. In an open
system requests arrive at a rate that does not depend on completions. Schroeder et al.\
showed that the two models behave very differently and that the model must match the system being
evaluated~\cite{Schroeder2006}.

A closed-loop generator used to measure latency suffers from coordinated omission: when the system
stalls, the generator stops sending, so it records one slow request instead of all the requests
that real users would have issued during the stall, and the reported tail is too
optimistic~\cite{Tene2015}. Constant-rate generators avoid this. wrk2 measures latency from the time each request was
scheduled~\cite{wrk2}. Vegeta, used in this study, issues requests at a constant rate regardless of response
times and measures latency from the actual send~\cite{Vegeta}. Load testers have further pitfalls that distort
tail measurements~\cite{Zhang2016Treadmill,Kogias2019}. A closed loop is the natural way to measure throughput as
a function of concurrency; its latency percentiles are not used in this study.

\subsection{The Universal Scalability Law}
\label{sec:bg-usl}

Gunther's Universal Scalability Law (USL) models the throughput of a closed system with $N$
concurrent clients as
\begin{equation}
  X(N) = \frac{\lambda N}{1 + \alpha (N-1) + \beta N (N-1)},
  \label{eq:usl}
\end{equation}
where $\lambda$ is the throughput with one client, $\alpha \ge 0$ is the contention coefficient and
$\beta \ge 0$ the coherency coefficient~\cite{Gunther2007}. Contention describes work that
serialises on a shared resource, such as a lock or a single saturated CPU, and on its own
($\beta = 0$) makes throughput approach $\lambda/\alpha$ as $N$ grows. Coherency describes the cost
of keeping shared state consistent between pairs of clients; it grows with $N(N-1)$ and, when
$\beta > 0$, makes throughput peak at
\begin{equation}
  N^{*} = \sqrt{(1-\alpha)/\beta}
  \label{eq:nstar}
\end{equation}
and decline beyond it. With $\alpha = \beta = 0$ scaling is linear. We write $\alpha$ and $\beta$
for the coefficients that Gunther denotes $\sigma$ and $\kappa$~\cite{Gunther2007}.

Because Equation~\eqref{eq:usl}
describes a closed system in which $N$ is the number of concurrent clients, this study fits it only
to closed-loop measurements taken at controlled concurrency levels. Open-loop measurements, where
concurrency is an outcome rather than an input, are used for latency.
